% ======================================================================
% SÉLECTION DE CONTENUS — Candidature type "Ingénierie Système"
% ======================================================================
% Ce fichier NE CONTIENT AUCUNE MISE EN PAGE. Il choisit juste, parmi les
% macros de cv_banque_contenus.tex, lesquelles afficher et dans quel ordre,
% pour CETTE candidature.
%
% Pour une autre offre : duplique ce fichier (ex: selection_thales.tex),
% modifie les 3 macros ci-dessous, puis change la ligne \input{...} dans
% cv_lucas_vacquie.tex pour pointer vers le nouveau fichier. Le layout
% principal n'a jamais besoin d'être touché.
%
% Rappel des blocs disponibles dans la banque :
%   Expérience : \expMicrotec \expPliageMetal \expGabarit
%   Associatif : \assoHandball \assoBenevole \assoGabarit
%   Projets    : \projBrasRobotique \projBuckConverter \projSatellite
%                \projBBEight \projQueueIA \projSoccerKit \projAppMobileAsso
%                \projNuitICAM \projExperiment \projGabarit
%
% Convention : chaque macro de sélection suppose qu'une coordonnée de départ
% (exp_pos / asso_pos / proj_pos) a déjà été positionnée par le CV principal
% juste avant l'appel (sous le titre de la rubrique).
% ======================================================================

% --- EXPÉRIENCE ---
\newcommand{\SelectionExperience}{%
    \expMicrotec{exp_pos}
    \chainNext{exp_pos}{exp_last}{\gapTextToSub}
    \expPliageMetal{exp_pos}
}

% --- ENGAGEMENT ASSOCIATIF ---
\newcommand{\SelectionAssociatif}{%
    \assoHandball{asso_pos}
    \chainNext{asso_pos}{asso_last}{\gapTextToSub}
    \assoBenevole{asso_pos}
}

% --- PROJETS RÉALISÉS ---
\newcommand{\SelectionProjets}{%
    \projBrasRobotique{proj_pos}
    \chainNext{proj_pos}{proj_last}{0.05cm}
    \projNuitICAM{proj_pos}
    \chainNext{proj_pos}{proj_last}{0.05cm}
    \projSatellite{proj_pos}
    \chainNext{proj_pos}{proj_last}{0.05cm}
    \projBBEight{proj_pos}
}
